\chapter{Complexity — SAT To3 SAT}
%%% generated by the blueprint pipeline from TCSlib.Complexity.NPReductions.SATTo3SAT
\section{Overview}

This module formalises the classical polynomial-time reduction from SAT to 3-SAT
via a chain (Tseitin-style) encoding.  Given a CNF formula $f$ over variables $V$,
the function \texttt{SATTo3SAT.to3SAT} produces an equisatisfiable 3-CNF formula over an
extended variable type \texttt{SATTo3SAT.AuxVar} $V$; the main result
\texttt{SATTo3SAT.SAT\_to\_3SAT\_equivalence} states that $f$ is satisfiable if and only if
its 3-CNF encoding is satisfiable.

\section{Declarations}

\begin{definition}[Clause]
\lean{SATTo3SAT.Clause}
\label{SATTo3SAT.Clause}
\leanok
A clause over variable type $V$ is a disjunction represented as a list of literals,
i.e.\ $\mathtt{Clause}\,V = \mathtt{List}\,(\mathtt{Literal}\,V)$.
\end{definition}

\begin{definition}[CNF formula]
\lean{SATTo3SAT.CNFFormula}
\label{SATTo3SAT.CNFFormula}
\leanok
\uses{SATTo3SAT.Clause}
A CNF formula over $V$ is a conjunction of clauses, represented as a list of clauses:
$\mathtt{CNFFormula}\,V = \mathtt{List}\,(\mathtt{Clause}\,V)$.
\end{definition}

\begin{definition}[Literal evaluation]
\lean{SATTo3SAT.evalLiteral}
\label{SATTo3SAT.evalLiteral}
\leanok
Given an assignment $\alpha : V \to \mathrm{Prop}$, \texttt{SATTo3SAT.evalLiteral} maps
a literal to its truth value: a positive literal $\mathtt{pos}\,v$ is true iff $\alpha(v)$
holds, and a negative literal $\mathtt{neg}\,v$ is true iff $\alpha(v)$ does not hold.
\end{definition}

\begin{definition}[Clause satisfaction]
\lean{SATTo3SAT.clauseSatisfied}
\label{SATTo3SAT.clauseSatisfied}
\leanok
\uses{SATTo3SAT.Clause, SATTo3SAT.evalLiteral}
A clause $c$ is satisfied by assignment $\alpha$ if at least one literal in $c$ evaluates
to true under $\alpha$, i.e.\ $\exists\, l \in c,\; \texttt{SATTo3SAT.evalLiteral}\,\alpha\,l$.
\end{definition}

\begin{definition}[Formula satisfaction]
\lean{SATTo3SAT.formulaSatisfied}
\label{SATTo3SAT.formulaSatisfied}
\leanok
\uses{SATTo3SAT.CNFFormula, SATTo3SAT.clauseSatisfied}
A CNF formula $f$ is satisfied by assignment $\alpha$ if every clause in $f$ is satisfied
by $\alpha$.
\end{definition}

\begin{definition}[Satisfiability of a CNF formula]
\lean{SATTo3SAT.isSatisfiable}
\label{SATTo3SAT.isSatisfiable}
\leanok
\uses{SATTo3SAT.CNFFormula, SATTo3SAT.formulaSatisfied}
A CNF formula $f$ is satisfiable if there exists an assignment $\alpha$ such that
$\texttt{SATTo3SAT.formulaSatisfied}\,\alpha\,f$ holds.
\end{definition}

\begin{definition}[3-clause]
\lean{SATTo3SAT.Clause3}
\label{SATTo3SAT.Clause3}
\leanok
A 3-clause over $V$ is a structure carrying exactly three literals $l_1, l_2, l_3 :
\mathtt{Literal}\,V$; a 3-clause is satisfied when at least one of the three literals is true.
\end{definition}

\begin{definition}[3-CNF formula]
\lean{SATTo3SAT.Formula3}
\label{SATTo3SAT.Formula3}
\leanok
\uses{SATTo3SAT.Clause3}
A 3-CNF formula over $V$ is a list of 3-clauses:
$\mathtt{Formula3}\,V = \mathtt{List}\,(\mathtt{Clause3}\,V)$.
\end{definition}

\begin{definition}[3-clause satisfaction]
\lean{SATTo3SAT.clause3Satisfied}
\label{SATTo3SAT.clause3Satisfied}
\leanok
\uses{SATTo3SAT.Clause3, SATTo3SAT.evalLiteral}
A 3-clause $c$ is satisfied by assignment $\alpha$ if at least one of its three literals
$c.l_1$, $c.l_2$, $c.l_3$ evaluates to true under $\alpha$.
\end{definition}

\begin{definition}[3-CNF formula satisfaction]
\lean{SATTo3SAT.formula3Satisfied}
\label{SATTo3SAT.formula3Satisfied}
\leanok
\uses{SATTo3SAT.Formula3, SATTo3SAT.clause3Satisfied}
A 3-CNF formula $f$ is satisfied by assignment $\alpha$ if every 3-clause in $f$ is
satisfied by $\alpha$.
\end{definition}

\begin{definition}[3-SAT satisfiability]
\lean{SATTo3SAT.is3Satisfiable}
\label{SATTo3SAT.is3Satisfiable}
\leanok
\uses{SATTo3SAT.Formula3, SATTo3SAT.formula3Satisfied}
A 3-CNF formula $f$ is 3-satisfiable if there exists an assignment $\alpha$ such that
$\texttt{SATTo3SAT.formula3Satisfied}\,\alpha\,f$ holds.
\end{definition}

\begin{definition}[Auxiliary variable type]
\lean{SATTo3SAT.AuxVar}
\label{SATTo3SAT.AuxVar}
\leanok
The inductive type $\mathtt{AuxVar}\,V$ extends a variable type $V$ with two constructors:
$\mathtt{orig}\,v$ wraps an original variable $v : V$, and $\mathtt{extra}\,i\,j$
introduces a fresh auxiliary variable $y_{i,j}$ used as a chain link when encoding clause
$i$ of length greater than three.
\end{definition}

\begin{definition}[Lift a literal to the extended type]
\lean{SATTo3SAT.liftLit}
\label{SATTo3SAT.liftLit}
\leanok
\uses{SATTo3SAT.AuxVar}
\texttt{SATTo3SAT.liftLit} maps a literal over $V$ to the corresponding literal over
$\mathtt{AuxVar}\,V$ by wrapping each variable with $\mathtt{orig}$, preserving polarity.
\end{definition}

\begin{definition}[Build chain clauses]
\lean{SATTo3SAT.buildChain}
\label{SATTo3SAT.buildChain}
\leanok
\uses{SATTo3SAT.AuxVar, SATTo3SAT.Clause3}
Given a clause index $c\_idx$, a lifting map $ml$, a suffix list of literals $\mathit{lits}$,
and a chain counter $j \ge 1$, \texttt{SATTo3SAT.buildChain} produces the middle and final
3-clauses of the chain encoding.  Each middle step emits $\langle \lnot y_{j-1},\, l_i,\, y_j\rangle$;
the final step emits $\langle \lnot y_{j-1},\, l_{n-1},\, l_n\rangle$ when only two literals remain.
\end{definition}

\begin{definition}[Encode a single clause as 3-clauses]
\lean{SATTo3SAT.transformClause}
\label{SATTo3SAT.transformClause}
\leanok
\uses{SATTo3SAT.AuxVar, SATTo3SAT.Clause3, SATTo3SAT.buildChain}
\texttt{SATTo3SAT.transformClause} maps a single clause (at index $c\_idx$) to a list of
3-clauses: the empty clause becomes two contradictory 3-clauses (unsatisfiable), a
1-literal clause becomes $\langle l_1,l_1,l_1\rangle$, a 2-literal clause becomes
$\langle l_1,l_2,l_2\rangle$, a 3-literal clause is copied as-is, and any longer clause
$l_1 :: l_2 :: \mathit{rest}$ becomes $\langle l_1, l_2, y_0\rangle$ followed by
$\texttt{SATTo3SAT.buildChain}$ on $\mathit{rest}$.
\end{definition}

\begin{definition}[Recursive SAT-to-3SAT worker]
\lean{SATTo3SAT.to3SATAux}
\label{SATTo3SAT.to3SATAux}
\leanok
\uses{SATTo3SAT.AuxVar, SATTo3SAT.Clause, SATTo3SAT.Clause3, SATTo3SAT.transformClause}
The auxiliary recursive function \texttt{SATTo3SAT.to3SATAux} iterates
\texttt{SATTo3SAT.transformClause} over a list of clauses, threading a clause index
counter to ensure each clause uses distinct auxiliary variables.
\end{definition}

\begin{definition}[SAT to 3-SAT encoding]
\lean{SATTo3SAT.to3SAT}
\label{SATTo3SAT.to3SAT}
\leanok
\uses{SATTo3SAT.AuxVar, SATTo3SAT.CNFFormula, SATTo3SAT.Formula3, SATTo3SAT.to3SATAux}
\texttt{SATTo3SAT.to3SAT} converts a CNF formula $f$ over $V$ into a 3-CNF formula over
$\mathtt{AuxVar}\,V$ by calling \texttt{SATTo3SAT.to3SATAux} starting at index $0$.
\end{definition}

\begin{definition}[Auxiliary variable valuation]
\lean{SATTo3SAT.extraVal}
\label{SATTo3SAT.extraVal}
\leanok
\uses{SATTo3SAT.evalLiteral}
For a list of literals $\mathit{lits}$ and index $j$, $\texttt{SATTo3SAT.extraVal}\,\alpha\,\mathit{lits}\,j$
holds iff every literal in $\mathit{lits}.\mathtt{take}(j + 2)$ is false under $\alpha$.
This is the intended truth value of the auxiliary chain variable $y_j$ in the encoding
of $\mathit{lits}$.
\end{definition}

\begin{lemma}[Auxiliary variable valuation as an all-false prefix]
\lean{SATTo3SAT.extraVal_iff_prefix_all_false}
\label{SATTo3SAT.extraVal_iff_prefix_all_false}
\leanok
\uses{SATTo3SAT.evalLiteral, SATTo3SAT.extraVal}
\difficulty{1}
Fix an assignment $\alpha$ and a list of literals $\mathit{lits}$, and let $j$ be an
index with $j + 2 \le |\mathit{lits}|$. Then the auxiliary variable valuation associated
with $\alpha$, $\mathit{lits}$, and $j$ holds if and only if every one of the first $j +
2$ literals of $\mathit{lits}$ evaluates to false under $\alpha$.
\end{lemma}

\begin{definition}[Global assignment lift]
\lean{SATTo3SAT.globalAssignment}
\label{SATTo3SAT.globalAssignment}
\leanok
\uses{SATTo3SAT.AuxVar, SATTo3SAT.CNFFormula, SATTo3SAT.extraVal}
Given $\alpha : \mathtt{Assignment}\,V$ and a formula $f$, \texttt{SATTo3SAT.globalAssignment}
$\alpha\,f$ is the assignment on $\mathtt{AuxVar}\,V$ that maps each original variable
$\mathtt{orig}\,v$ to $\alpha(v)$ and each auxiliary variable $\mathtt{extra}\,i\,j$ to
$\texttt{SATTo3SAT.extraVal}\,\alpha\,(f.\mathtt{get}\,i)\,j$, i.e.\ the chain invariant
``the first $j+2$ literals of clause $i$ are all false.''
\end{definition}

\begin{lemma}[Membership in the clause-by-clause 3-SAT encoding]
\lean{SATTo3SAT.mem_to3SATAux_iff}
\label{SATTo3SAT.mem_to3SATAux_iff}
\leanok
\uses{SATTo3SAT.AuxVar, SATTo3SAT.Clause, SATTo3SAT.Clause3, SATTo3SAT.to3SATAux, SATTo3SAT.transformClause}
\difficulty{4}
Let $\mathit{cs} = (c_0,\dots,c_{m-1})$ be a list of clauses over a variable type $V$,
let $\mathit{idx} \in \bbn$ be a starting index, and let $c_3$ be a 3-clause over the
auxiliary variable type $\mathtt{AuxVar}\,V$. Encode $\mathit{cs}$ into 3-clauses by
transforming each clause $c_i$ at the shifted index $\mathit{idx}+i$ and concatenating
the outputs in order of increasing $i$. Then $c_3$ occurs in the resulting encoding of
$\mathit{cs}$ if and only if there is a position $i < m$ for which $c_3$ occurs among
the 3-clauses produced by transforming the clause $c_i$ at index $\mathit{idx}+i$.
\end{lemma}

\begin{lemma}[Membership in the SAT-to-3-SAT encoding]
\lean{SATTo3SAT.mem_to3SAT_iff}
\label{SATTo3SAT.mem_to3SAT_iff}
\leanok
\uses{SATTo3SAT.AuxVar, SATTo3SAT.CNFFormula, SATTo3SAT.Clause3, SATTo3SAT.mem_to3SATAux_iff, SATTo3SAT.to3SAT, SATTo3SAT.transformClause}
\difficulty{2}
Let $f = (C_0, \dots, C_{m-1})$ be a CNF formula over a variable type $V$, and write
$\mathrm{Aux}(V)$ for the auxiliary variable type built from $V$. For each index $i$
with $0 \le i < m$, let $T_i$ be the list of $3$-clauses obtained by encoding the single
clause $C_i$ at index $i$, and take the SAT-to-3-SAT encoding of $f$ to be the $3$-CNF
formula over $\mathrm{Aux}(V)$ formed by concatenating $T_0, T_1, \dots, T_{m-1}$ in
order. Then a $3$-clause $c$ over $\mathrm{Aux}(V)$ occurs in the SAT-to-3-SAT encoding
of $f$ if and only if there exists an index $i < m$ such that $c$ occurs in $T_i$.
\end{lemma}

\begin{lemma}[Extending a list prefix past a dropped head]
\lean{SATTo3SAT.List.take_succ_of_drop_cons}
\label{SATTo3SAT.List.take_succ_of_drop_cons}
\leanok
\difficulty{4}
Let $l$ be a list over a type $\alpha$, let $n$ be a natural number, and suppose that
dropping the first $n$ entries of $l$ leaves a nonempty list whose head is $a$ (so that
$l$ has more than $n$ entries and its $(n+1)$-st entry is $a$). Then the prefix
consisting of the first $n+1$ entries of $l$ equals the prefix consisting of the first
$n$ entries with $a$ appended: taking $n+1$ entries of $l$ gives (taking $n$ entries of
$l$) followed by $[a]$.
\end{lemma}

\begin{lemma}[Chain clauses satisfied by the completed assignment]
\lean{SATTo3SAT.buildChain_all_satisfied}
\label{SATTo3SAT.buildChain_all_satisfied}
\leanok
\uses{SATTo3SAT.AuxVar, SATTo3SAT.List.take_succ_of_drop_cons, SATTo3SAT.buildChain, SATTo3SAT.clause3Satisfied, SATTo3SAT.evalLiteral, SATTo3SAT.extraVal}
\difficulty{6}
Fix a clause index $i$ and a clause presented as a list of literals over a variable type
$V$, and let $\alpha : V \to \mathrm{Prop}$ be an assignment that makes at least one
literal of the clause true. Let $ml$ be a map lifting each literal over $V$ to a literal
over the extended type $\mathtt{AuxVar}\,V$, and let $\alpha_3$ be an assignment on
$\mathtt{AuxVar}\,V$ such that: (i) for every $k$ the auxiliary variable $y_{i,k}$ is
true under $\alpha_3$ exactly when all of the first $k+2$ literals of the clause are
false under $\alpha$; and (ii) for every literal $\ell$, the lifted literal $ml(\ell)$
is true under $\alpha_3$ if and only if $\ell$ is true under $\alpha$. Then for every
chain counter $j \ge 1$, letting $\mathit{lits}$ be the suffix of the clause obtained by
dropping its first $j+1$ literals, every 3-clause in the chain built for clause $i$ from
$\mathit{lits}$ starting at counter $j$ is satisfied by $\alpha_3$.
\end{lemma}

\begin{definition}[Local literal predicate]
\lean{SATTo3SAT.IsLocalVar}
\label{SATTo3SAT.IsLocalVar}
\leanok
\uses{SATTo3SAT.AuxVar}
A literal $l$ in $\mathtt{Literal}\,(\mathtt{AuxVar}\,V)$ is \emph{local to index $i$} if
every auxiliary variable appearing in $l$ carries index $i$.  Literals over original
variables are always local.
\end{definition}

\begin{lemma}[Local and global assignments agree on local literals]
\lean{SATTo3SAT.eval_local_eq_global}
\label{SATTo3SAT.eval_local_eq_global}
\leanok
\uses{SATTo3SAT.AuxVar, SATTo3SAT.CNFFormula, SATTo3SAT.Clause, SATTo3SAT.IsLocalVar, SATTo3SAT.evalLiteral, SATTo3SAT.extraVal, SATTo3SAT.globalAssignment}
\difficulty{3}
Let $\alpha$ be a truth assignment of the variables $V$, and let $f$ be a CNF formula
over $V$ with clauses $C_0, \dots, C_{m-1}$. Fix an index $i$ with $0 \le i < m$ and
write $C_i$ for the $i$-th clause. Let $l$ be a literal over the auxiliary variable type
$\mathtt{AuxVar}\,V$ that is local to $i$, meaning that every auxiliary variable
$y_{i',j}$ occurring in $l$ satisfies $i' = i$. Write $\hat\alpha$ for the global
assignment determined by $\alpha$ and $f$, which sends each original variable to its
value under $\alpha$ and each auxiliary variable $y_{i',j}$ to the truth of the
assertion that the first $j+2$ literals of clause $C_{i'}$ are all false under $\alpha$;
and write $\alpha_i$ for the corresponding local assignment, built from clause $C_i$
alone. Then $l$ takes the same truth value under $\hat\alpha$ as under $\alpha_i$.
\end{lemma}

\begin{lemma}[Chain-encoding clauses stay local to their index]
\lean{SATTo3SAT.buildChain_isLocal}
\label{SATTo3SAT.buildChain_isLocal}
\leanok
\uses{SATTo3SAT.AuxVar, SATTo3SAT.Clause3, SATTo3SAT.IsLocalVar, SATTo3SAT.buildChain}
\difficulty{4}
Let $i$ be a natural number and let $ml$ be a lifting map sending each literal over a
variable type $V$ to a literal over $\mathtt{AuxVar}\,V$, and suppose that $ml(l)$ is
local to index $i$ for every literal $l$. Then, for every list of literals
$\mathit{lits}$ over $V$ and every chain counter $j$, each $3$-clause occurring in the
chain of clauses built from $\mathit{lits}$ (with clause index $i$, lifting map $ml$,
and counter $j$) has all three of its literals local to index $i$.
\end{lemma}

\begin{lemma}[Locality of the 3-clause encoding of a single clause]
\lean{SATTo3SAT.transformClause_isLocal}
\label{SATTo3SAT.transformClause_isLocal}
\leanok
\uses{SATTo3SAT.AuxVar, SATTo3SAT.Clause, SATTo3SAT.Clause3, SATTo3SAT.IsLocalVar, SATTo3SAT.buildChain_isLocal, SATTo3SAT.transformClause}
\difficulty{4}
Fix a clause index $i$ and let $C$ be a clause over a variable type $V$. If a 3-clause
occurs in the list produced by encoding $C$ at index $i$, then each of its three
literals is local to $i$; that is, every auxiliary variable appearing in the 3-clause
carries the index $i$, while literals over original variables of $V$ are local
automatically.
\end{lemma}

\begin{lemma}[Global and local assignments agree on an encoded clause]
\lean{SATTo3SAT.global_matches_local}
\label{SATTo3SAT.global_matches_local}
\leanok
\uses{SATTo3SAT.AuxVar, SATTo3SAT.CNFFormula, SATTo3SAT.Clause, SATTo3SAT.Clause3, SATTo3SAT.clause3Satisfied, SATTo3SAT.eval_local_eq_global, SATTo3SAT.extraVal, SATTo3SAT.globalAssignment, SATTo3SAT.transformClause, SATTo3SAT.transformClause_isLocal}
\difficulty{2}
Fix an assignment $\alpha$ of truth values to the variables of a type $V$, a CNF formula
$f$, and an index $i$ with $i < \mathrm{length}(f)$; let $C$ be the $i$-th clause of
$f$, and let $c_3$ be any of the 3-clauses produced by encoding $C$ as clause number
$i$. Consider the global lifted assignment on $\mathtt{AuxVar}\,V$, which sends each
original variable $v$ to $\alpha(v)$ and each auxiliary chain variable $y_{i',j}$ to the
truth value of the assertion ``the first $j+2$ literals of the $i'$-th clause of $f$ are
all false under $\alpha$'', together with the local assignment that agrees with $\alpha$
on original variables and sends $y_{i,j}$ to the truth value of ``the first $j+2$
literals of $C$ are all false under $\alpha$''. Then $c_3$ is satisfied by the global
assignment if and only if it is satisfied by the local assignment.
\end{lemma}

\begin{lemma}[Extended assignment satisfies the 3-clause encoding of a satisfied clause]
\lean{SATTo3SAT.transformClause_satisfied}
\label{SATTo3SAT.transformClause_satisfied}
\leanok
\uses{SATTo3SAT.AuxVar, SATTo3SAT.Clause, SATTo3SAT.buildChain_all_satisfied, SATTo3SAT.clause3Satisfied, SATTo3SAT.evalLiteral, SATTo3SAT.extraVal, SATTo3SAT.transformClause}
\difficulty{5}
Let $\alpha$ be a truth assignment on a variable set $V$, and let $C$ be a clause over
$V$ — a finite list of literals read as their disjunction — carrying the index $i$.
Assume $\alpha$ satisfies $C$, meaning at least one literal of $C$ is true under
$\alpha$. Extend $\alpha$ to an assignment $\alpha_3$ on the enlarged variable set
consisting of the original variables of $V$ together with the auxiliary chain variables
$y_{i,j}$: leave each original variable at its value under $\alpha$, and set $y_{i,j}$
to true exactly when the first $j+2$ literals of $C$ are all false under $\alpha$. Then
every $3$-clause in the list obtained by encoding $C$ at index $i$ as $3$-clauses is
satisfied by $\alpha_3$.
\end{lemma}

\begin{theorem}[Completeness of the SAT to 3-SAT encoding]
\lean{SATTo3SAT.SAT_to_3SAT_completeness}
\label{SATTo3SAT.SAT_to_3SAT_completeness}
\leanok
\uses{SATTo3SAT.CNFFormula, SATTo3SAT.clauseSatisfied, SATTo3SAT.globalAssignment, SATTo3SAT.global_matches_local, SATTo3SAT.is3Satisfiable, SATTo3SAT.isSatisfiable, SATTo3SAT.mem_to3SAT_iff, SATTo3SAT.to3SAT, SATTo3SAT.transformClause_satisfied}
\difficulty{3}
Let $f$ be a CNF formula over a variable type $V$, and let $\mathrm{to3SAT}(f)$ be the
3-CNF formula produced by the standard clause-splitting transformation, whose variables
are those of $V$ together with fresh auxiliary chain variables introduced to break each
long clause into a chain of three-literal clauses. If $f$ is satisfiable — that is, some
assignment to $V$ makes every clause of $f$ true — then $\mathrm{to3SAT}(f)$ is
3-satisfiable: there is an assignment to its variables under which every 3-clause is
true.
\end{theorem}

\begin{lemma}[Contradiction from a satisfied chain of all-false literals]
\lean{SATTo3SAT.buildChain_forced_false}
\label{SATTo3SAT.buildChain_forced_false}
\leanok
\uses{SATTo3SAT.AuxVar, SATTo3SAT.buildChain, SATTo3SAT.clause3Satisfied, SATTo3SAT.evalLiteral}
\difficulty{5}
Fix a clause index $i$ and an assignment $\alpha$ of truth values to the auxiliary
variables. Let $\ell \mapsto \widehat{\ell}$ be a map lifting each literal over the base
variable type to a literal over the auxiliary-variable type, let $j \ge 1$, and let $L$
be a list of at least two literals over the base type such that the lifted literal
$\widehat{\ell}$ evaluates to false under $\alpha$ for every $\ell$ in $L$. If every
3-clause of the chain encoding built from clause index $i$, the lifting $\ell \mapsto
\widehat{\ell}$, the list $L$, and the counter $j$ is satisfied by $\alpha$, then
$\alpha$ cannot also make the entry variable $y_{i,\,j-1}$ true; these hypotheses taken
together are contradictory.
\end{lemma}

\begin{lemma}[Soundness of the clause-to-3-clause encoding]
\lean{SATTo3SAT.transformClause_soundness}
\label{SATTo3SAT.transformClause_soundness}
\leanok
\uses{SATTo3SAT.AuxVar, SATTo3SAT.Clause, SATTo3SAT.buildChain_forced_false, SATTo3SAT.clause3Satisfied, SATTo3SAT.evalLiteral, SATTo3SAT.transformClause}
\difficulty{5}
Fix a variable type $V$, a clause index $c \in \bbn$, and a clause $C$ over $V$ — a
finite list of literals read as their disjunction. Let $\alpha$ be a truth assignment to
the auxiliary variables (the original variables of $V$ together with the fresh chain
variables), and suppose that every $3$-clause produced by encoding $C$ at index $c$ is
satisfied by $\alpha$. Then $C$ has at least one literal that is true under the
assignment obtained by reading each original variable $v$ through its copy among the
auxiliary variables, i.e.\ $v \mapsto \alpha(\mathrm{orig}\,v)$.
\end{lemma}

\begin{theorem}[Soundness of the SAT-to-3-SAT encoding]
\lean{SATTo3SAT.SAT_to_3SAT_soundness}
\label{SATTo3SAT.SAT_to_3SAT_soundness}
\leanok
\uses{SATTo3SAT.CNFFormula, SATTo3SAT.is3Satisfiable, SATTo3SAT.isSatisfiable, SATTo3SAT.mem_to3SAT_iff, SATTo3SAT.to3SAT, SATTo3SAT.transformClause_soundness}
\difficulty{3}
Let $f$ be a CNF formula over a variable type $V$, and let its SAT-to-3-SAT encoding be
the 3-CNF formula over the extended variable type (the original variables together with
the fresh auxiliary chain variables). If this encoded 3-CNF formula is 3-satisfiable —
that is, some assignment of the extended variables satisfies all of its 3-clauses — then
$f$ itself is satisfiable.
\end{theorem}

\begin{theorem}[SAT-to-3-SAT preserves satisfiability]
\lean{SATTo3SAT.SAT_to_3SAT_equivalence}
\label{SATTo3SAT.SAT_to_3SAT_equivalence}
\leanok
\uses{SATTo3SAT.CNFFormula, SATTo3SAT.SAT_to_3SAT_completeness, SATTo3SAT.SAT_to_3SAT_soundness, SATTo3SAT.is3Satisfiable, SATTo3SAT.isSatisfiable, SATTo3SAT.to3SAT}
\difficulty{1}
Let $f$ be a CNF formula over a variable type $V$, and let $\tau(f)$ denote its $3$-CNF
encoding: the formula obtained by replacing each clause of $f$ with an equisatisfiable
chain of $3$-clauses over an enlarged variable type that adjoins to $V$ a supply of
fresh auxiliary variables (a clause of more than three literals is split into linked
triples using these auxiliary variables, while shorter clauses are padded to three
literals). Then $f$ is satisfiable if and only if $\tau(f)$ is $3$-satisfiable.
\end{theorem}
